% Альтернативная версия tex/adp_matrix_statistics.tex.
% Файл использует обозначения и макросы из tex/manifold_v2.tex
% и должен подключаться после него.

\section{Батчевое вычисление статистик ADP}
\label{sec:adp-batched-statistics}

В этом разделе описана реализация из
\texttt{ADP/ADP\_Statistic.py}:
\begin{center}
    \texttt{calculate\_statistics} (CPU),\qquad
    \texttt{calculate\_statistics\_gpu} (GPU).
\end{center}
Используются обозначения
раздела~\ref{Ssingind}: наблюдения
$\Xv_i\in\R^{\dimd}$, отклики $Y_i\in\R$, центры
$\xv_j$, веса
\begin{equation}
    \weight_{ij}
    =
    \Kern\!\left(\left\|\TEDR(\Xv_i-\xv_j)\right\|^2\right),
    \qquad i=1,\ldots,n,
    \quad j\in\JJJ,
\label{eq:batch-adp-weights}
\end{equation}
и направления $\sph_{jr}\in\Sph_j$, $r=1,\ldots,\nSph$.
В коде массив направлений имеет форму
$\nJJJ\times\nSph\times\dimd$.

В отличие от определения через исходный центр $\xv_j$ в
\eqref{ytd7e7ydwuhd5eythdm}, численная реализация центрирует каждую
окрестность относительно её взвешенного среднего
\begin{equation}
    \Mv_j
    =
    \frac{1}{\mu_j}\sum_{i=1}^n \weight_{ij}\Xv_i,
    \qquad
    \mu_j=\sum_{i=1}^n\weight_{ij}>0.
\label{eq:batch-adp-local-mean}
\end{equation}

Поэтому вычисляемые статистики имеют вид
\begin{align}
    \Ima_{j,\sph_{jr}}
    &=
    \sum_{i=1}^n
    \weight_{ij}(Y_i-\bar Y_j)
    \langle \Xv_i-\Mv_j,\sph_{jr}\rangle,
\label{eq:batch-adp-I}
    \\
    \Uv_{j,\sph_{jr}}
    &=
    \sum_{i=1}^n
    \weight_{ij}(\Xv_i-\Mv_j)
    \langle \Xv_i-\Mv_j,\sph_{jr}\rangle,
\label{eq:batch-adp-U}
\end{align}
где
$\bar Y_j=\mu_j^{-1}\sum_i\weight_{ij}Y_i$.
В точной арифметике центрирование $Y_i$ в
\eqref{eq:batch-adp-I} не меняет результат, поскольку
\begin{equation}
    \sum_{i=1}^n
    \weight_{ij}\langle\Xv_i-\Mv_j,\sph_{jr}\rangle=0.
\label{eq:batch-adp-zero-moment}
\end{equation}

\subsection{Разбиение по центрам}

Пусть
$\mathcal B=\{j_0,\ldots,j_0+B-1\}\subset\JJJ$
--- очередной батч центров. Для него вводятся
\begin{equation}
    W_{\mathcal B}
    =
    (\weight_{ij})_{j\in\mathcal B,\,i=1}^n
    \in\R^{B\times n},
    \qquad
    A_{\mathcal B}
    =
    \operatorname{diag}(\mu_{\mathcal B})^{-1}W_{\mathcal B}.
\label{eq:batch-adp-normalized-weights}
\end{equation}
Строки $A_{\mathcal B}$ суммируются в единицу.
Полная матрица $W\in\R^{\nJJJ\times n}$, если она передана явно,
разрезается на такие блоки параметром \texttt{batch\_size}. Вместо неё
можно передать итератор уже сформированных блоков. Блоки обязаны без
пропусков и перестановок покрывать центры от $0$ до $\nJJJ-1$; размер
блока итератора задаётся его производителем и повторно не разрезается.

Для уменьшения ошибки при больших общих сдвигах сначала выполняется
глобальное центрирование
\begin{equation}
    \bar\Xv=\frac1n\sum_{i=1}^n\Xv_i,
    \qquad
    \Xv_i^c=\Xv_i-\bar\Xv,
    \qquad
    \Mv_j^c=\sum_{i=1}^na_{ji}\Xv_i^c.
\label{eq:batch-adp-global-centering}
\end{equation}
После вычислений возвращается
$\Mv_j=\Mv_j^c+\bar\Xv$.

Для всех $j\in\mathcal B$ и всех направлений одновременно строится
тензор проекций
\begin{equation}
    Q_{jri}
    =
    \langle\Xv_i^c-\Mv_j^c,\sph_{jr}\rangle,
    \qquad
    Q\in\R^{B\times\nSph\times n}.
\label{eq:batch-adp-Q}
\end{equation}
Из-за округления его нормированное взвешенное среднее
\begin{equation}
    e_{jr}=\sum_{i=1}^na_{ji}Q_{jri}
\label{eq:batch-adp-residual}
\end{equation}
может отличаться от нуля. Реализация явно вычитает этот остаток:
\begin{equation}
    H_{jri}=a_{ji}(Q_{jri}-e_{jr}),
    \qquad
    s_{jr}=\sum_{i=1}^nH_{jri}.
\label{eq:batch-adp-H}
\end{equation}
Тогда батч статистик вычисляется матричными произведениями
\begin{align}
    \Ima_{j,\sph_{jr}}
    &=
    \mu_j\left(
        \sum_{i=1}^nH_{jri}Y_i-s_{jr}\bar Y_j
    \right),
\label{eq:batch-adp-I-stable}
    \\
    \Uv_{j,\sph_{jr}}
    &=
    \mu_j\left(
        \sum_{i=1}^nH_{jri}\Xv_i^c-s_{jr}\Mv_j^c
    \right).
\label{eq:batch-adp-U-stable}
\end{align}
Члены с $s_{jr}$ равны нулю в точной арифметике, но защищают результат
от умножения малого остатка на большой общий сдвиг данных.

\subsection{Плотная ветвь}

Пусть $K_{\mathcal B}$ --- максимальное число ненулевых весов в одной
строке текущего батча. Если
\begin{equation}
    4K_{\mathcal B}>n,
\label{eq:batch-adp-dense-dispatch}
\end{equation}
реализация оставляет блок весов плотным. В терминах массивов операции
имеют вид
\begin{align}
    \Mv_{\mathcal B}^c
    &=A_{\mathcal B}X^c,
    \\
    Q
    &=\operatorname{bmm}\!\left(
        \Phi_{\mathcal B},(X^c)^{\mathsf T}
      \right)
      -\left(
        \langle\Mv_j^c,\sph_{jr}\rangle
       \right)_{j,r,\mathrm{None}},
    \\
    H
    &=(Q-e_{:\,,:\,,\mathrm{None}})
      \odot A_{\mathcal B,\,:\,,\mathrm{None},:}.
\end{align}
Здесь произведения выполняются независимо для каждой пары
$(j,r)$ с помощью batched matrix multiplication и broadcasting.
Основной временный массив $Q$ имеет форму
$B\times\nSph\times n$; массив результата $\Uv$ имеет форму
$B\times\nSph\times\dimd$. Поэтому рабочая память плотной ветви равна
\begin{equation}
    O(B\nSph n+B\nSph\dimd+Bn+B\dimd),
\end{equation}
а основная арифметика ---
$O(B\nSph n\dimd)$.

\subsection{Локальная и разреженная ветви}

Если $4K_{\mathcal B}\leq n$, плотный блок сначала преобразуется в
списки локальных соседей. Для каждой строки сохраняются индексы
\begin{equation}
    N\in\mathbb N^{B\times K_{\mathcal B}}
\end{equation}
и соответствующие веса
$W^{\mathrm{loc}}\in\R^{B\times K_{\mathcal B}}$.
Короткие строки дополняются нулевыми весами. После gather-операций
\begin{equation}
    X_{\mathrm{loc}}=X^c[N]
    \in\R^{B\times K_{\mathcal B}\times\dimd},
    \qquad
    Y_{\mathrm{loc}}=Y[N]
    \in\R^{B\times K_{\mathcal B}}
\end{equation}
формулы \eqref{eq:batch-adp-Q}--\eqref{eq:batch-adp-U-stable}
применяются только к локальным наблюдениям.

Объект \texttt{SparseNeighborhoodBlock} уже хранит окрестности в
сжатой форме: границы строк, индексы соседей и веса граничных точек.
Внутри \texttt{ADP\_Statistic.py} он также переводится в описанные выше
прямоугольные массивы $N$ и $W^{\mathrm{loc}}$. Следовательно, текущая
реализация не выполняет CSR--матричных произведений непосредственно;
сжатое представление экономит память между вызовами, а вычислительное
ядро использует padded gather и batched matrix multiplication.

Для локальной ветви рабочая память и арифметика оцениваются как
\begin{equation}
    O(BK_{\mathcal B}\dimd+B\nSph K_{\mathcal B}
      +B\nSph\dimd)
    \quad\text{и}\quad
    O(B\nSph K_{\mathcal B}\dimd)
\end{equation}
соответственно. Эти оценки зависят от максимальной ширины строки
в батче из-за padding.

Переданный \texttt{SparseStatisticsCache} регистрирует повторное
использование тех же окрестностей между последовательными вызовами.
Он не заменяет вычисление статистик и не хранит их численные значения.

\subsection{Диагностики и возвращаемые величины}

Помимо $\Ima$ и $\Uv$, для каждого центра возвращаются масса $\mu_j$,
среднее $\Mv_j$ и эффективное число наблюдений
\begin{equation}
    n_{\mathrm{eff},j}
    =
    \frac{1}{\sum_{i=1}^na_{ji}^2}
    =
    \frac{\left(\sum_i\weight_{ij}\right)^2}
         {\sum_i\weight_{ij}^2}.
\label{eq:batch-adp-neff}
\end{equation}
До вычитания $e_{jr}$ вычисляется относительная ошибка нулевого момента
\begin{equation}
    \eta_{jr}
    =
    \begin{cases}
    \displaystyle
    \frac{|e_{jr}|}{\sum_i a_{ji}|Q_{jri}|},
        & \sum_i a_{ji}|Q_{jri}|\neq0,\\[3mm]
    0,  & \text{иначе}.
    \end{cases}
\label{eq:batch-adp-eta}
\end{equation}
Малое $n_{\mathrm{eff},j}$ означает концентрацию массы на малом числе
точек, а большое $\eta_{jr}$ указывает на потерю точности при
центрировании проекций.

Все блоки должны содержать конечные неотрицательные веса и иметь
положительную конечную массу. Входные массивы должны иметь формы
\begin{equation}
    X:n\times\dimd,
    \qquad
    Y:n,
    \qquad
    \Phi:\nJJJ\times\nSph\times\dimd.
\end{equation}
После прохода проверяется, что блоки покрыли ровно $\nJJJ$ центров.
Массивы результата имеют формы
\begin{align}
    \Ima&:\nJJJ\times\nSph,
    &
    \Uv&:\nJJJ\times\nSph\times\dimd,
    \\
    \mu,\ n_{\mathrm{eff}}&:\nJJJ,
    &
    \Mv&:\nJJJ\times\dimd,
    &
    \eta&:\nJJJ\times\nSph.
\end{align}
Полные массивы результата сохраняются, а временные $W_{\mathcal B}$,
$Q$ и $H$ существуют только внутри текущего батча.

\subsection{CPU и GPU}

Функция \texttt{calculate\_statistics} выполняет все операции в NumPy.
Функция \texttt{calculate\_statistics\_gpu} переносит $X$, $Y$,
$\Phi$ и текущий блок весов на CUDA-устройство и применяет те же
плотные или локальные формулы в CuPy. Автоматического перехода на CPU
при отсутствии CuPy или CUDA нет.

В GPU-варианте тензоры $\Ima$ и $\Uv$ остаются на устройстве, поскольку
они непосредственно передаются GPU-реализации решателя. Масса $\mu$,
локальное среднее $\Mv$, $n_{\mathrm{eff}}$ и $\eta$ возвращаются как
NumPy-массивы. Все $\nSph$ направлений центра обрабатываются вместе;
параметр \texttt{batch\_size} ограничивает только число центров в
плотном блоке.
